%!TEX root = ../../thesis.tex

\section{Comparing the Two Formulations}\label{sec:results-head-to-head}

% Answers RQ2 (\cref{sec:introduction-rqs}).

% TODO: THE PREMISE OF THIS SECTION CHANGED --- rewrite it, do not patch it. It was
% written for a controlled comparison in which cJSON ran under both formulations. No
% target was run under both: the input formulation carries sequence alone
% (\cref{sec:targets-input-actions}) and the corpus formulation carries the three
% parsers (\cref{sec:targets-corpus-actions}). There is therefore NO paired
% experiment, and the section must not read as though there were.
%
% Decide which of these the section is, and commit:
%   (a) a design comparison rather than a measurement, arguing from the structure of
%       the two MDPs (\cref{chap:fuzzing-mdp}) and from what each one achieved on its
%       own targets, with the absence of a paired experiment stated as the first
%       sentence, or
%   (b) drop the section and answer RQ2 from \cref{chap:fuzzing-mdp} plus
%       \cref{sec:discussion-advantages}, which is the honest option if there is
%       nothing to report beyond the design argument.
% Do not leave a section whose title promises a comparison the data cannot support.
%
% If (a), the argument that IS available without a paired run:
%   - what each formulation can express. The input formulation builds one input left
%     to right and can represent sequential dependence between symbols directly; the
%     corpus formulation edits a persistent seed and can represent the classic fuzzer
%     loop but not a symbol's position in a protocol,
%   - what each costs per agent decision, and what an execution means in each. Under
%     the corpus formulation the place phase costs no execution at all, so half the
%     MDP steps are free; under the input formulation every step executes,
%   - why the input formulation was not carried to a real target
%     (\cref{sec:methodology-setup}). This is the decision RQ2 actually turned on, and
%     it should be argued rather than merely reported: an input-level environment must
%     re-drive the program per symbol, which for a batch parser means re-parsing a
%     growing prefix and paying for a harness that has no counterpart in how the
%     program is really used. The corpus formulation's persistent-mode loop, by
%     contrast, presumes exactly what a parser already is --- one self-contained
%     record per execution.
%
% Whatever the section becomes, name the confound that would have bounded the paired
% comparison anyway, because it is why the paired experiment is not merely missing but
% hard: the harness differs by necessity between the formulations, so an execution does
% not mean the same thing on both sides, and any coverage number compared across them
% compares two experimental apparatuses as much as two action spaces.
%
% Close with the transferable design lesson --- what the action-space choice costs and
% buys, and which formulation to reach for under which conditions. That conclusion
% feeds \cref{sec:discussion-advantages,sec:discussion-disadvantages}.
